\documentclass[12pt,a4paper]{article}
\usepackage[utf8]{inputenc}
\usepackage[T1]{fontenc}
\usepackage[french]{babel}
\usepackage{graphicx}
\usepackage{hyperref}
\usepackage{listings}
\usepackage{xcolor}
\usepackage{geometry}
\geometry{margin=2.5cm}

\title{\textbf{Présentation Technique \& Logique}\\ Moteur de Résolution (Solver) SBS/RS}
\author{Documentation Analytique}
\date{\today}

\begin{document}

\maketitle
\tableofcontents
\newpage

\section{Introduction}
Ce document synthétise l'architecture et la logique algorithmique du script \texttt{solver.py}. Ce module constitue l'intelligence décisionnelle du système : il prend en entrée un problème généré (un état initial désordonné et un état cible) et exécute des algorithmes pour rapprocher l'entrepôt de son état optimal.

L'enjeu central est de minimiser la distance totale parcourue par la navette, tout en évitant les blocages physiques (Deadlocks) causés par la contrainte de double profondeur.

\section{Concepts Techniques et Optimisations}

\subsection{Indexation Bidirectionnelle en O(1)}
Pour traiter des instances complexes rapidement, le \texttt{Solveur} n'utilise pas de boucles de recherche lourdes. Il maintient une indexation bidirectionnelle en mémoire vive :
\begin{itemize}
    \item \texttt{affectation} : Un dictionnaire associant \textit{Identifiant du Bac $\rightarrow$ Emplacement Physique}.
    \item \texttt{grille} : L'index inverse associant \textit{Emplacement Physique $\rightarrow$ Identifiant du Bac}.
\end{itemize}
Ces deux vues sont mises à jour instantanément à chaque déplacement (méthode \texttt{\_executer\_mouvement}), garantissant des performances d'accès de l'ordre de $O(1)$.

\subsection{Analyse de Graphe (Dépendance et Deadlocks)}
Le système évalue les interdépendances logiques entre les bacs. Un bac $A$ est dit dépendant d'un bac $B$ si le déplacement de $A$ vers sa cible finale est bloqué par la présence de $B$.

La méthode \texttt{\_suivre\_chaine()} parcourt ce graphe de dépendances. Si elle identifie une chaîne se terminant par un bac déplaçable, le mouvement est exécuté. Si elle détecte un \textbf{Cycle} ($A \rightarrow B \rightarrow C \rightarrow A$), cela signifie qu'il y a un deadlock : la seule solution est d'exécuter un mouvement temporaire, appelé \textit{Double Handling}, pour briser le cycle.

\subsection{Le "Double Handling" et la Mémoire Tabou}
Lorsqu'un bac bloque le système, la méthode \texttt{\_degager()} est appelée. Son rôle est de trouver la meilleure place vide temporaire. Pour éviter qu'un bac ne fasse des allers-retours infinis entre deux places, un mécanisme de \textbf{Liste Tabou} mémorise les dernières positions d'un bac et en interdit l'accès pendant les itérations suivantes.

\section{Explication Détaillée des 7 Heuristiques}
Le Solveur implémente 7 règles métier (heuristiques) différentes pour prioriser les bacs à déplacer et choisir où les stocker temporairement.

\subsection{H1 : Nearest Neighbor (Le Plus Proche Voisin)}
Il s'agit de la stratégie de base, focalisée sur la myopie spatiale. L'algorithme trie les bacs mal placés en fonction de leur proximité avec la position actuelle de la navette. S'il doit effectuer un Double Handling, il choisira la case vide la plus proche. Efficace en rapidité, mais génère souvent des sous-optimisations à long terme.

\subsection{H2 : Constrained Nearest Neighbor}
Une amélioration directe de H1 visant à limiter l'empreinte spatiale des opérations de dégagement. Lors d'un Double Handling, H2 oblige le bac bloquant à être stocké dans une case vide de la \textbf{même colonne}. S'il n'y a pas de place dans la colonne, l'heuristique recourt au comportement H1 en "fallback".

\subsection{H3 : Depth-First (Priorité au Remplissage Arrière)}
Cette heuristique s'attaque spécifiquement à la contrainte de la double profondeur. Au lieu de trier par distance stricte, elle priorise catégoriquement tous les bacs dont la cible finale se situe en \textbf{profondeur Arrière}. L'objectif logistique est de "boucher le fond" de l'entrepôt le plus vite possible pour éviter que la navette n'ait à bouger des bacs posés à l'avant (qui bloqueraient l'arrière).

\subsection{H4 : Tri Topologique sur Graphe}
Plutôt que de regarder l'espace physique, H4 analyse le graphe de dépendance. Elle attribue des scores pour classer les bacs. Sa priorité absolue est de s'attaquer aux \textbf{cycles} (deadlocks purs). Ensuite, elle trie selon la longueur des chaînes de blocage : en débloquant le nœud racine d'une longue chaîne, H4 espère créer un "effet domino" libérant de multiples bacs consécutifs.

\subsection{H5 : Look-Ahead Nearest Neighbor (Anticipation)}
Une variante intelligente de H1. Lorsqu'un Double Handling est nécessaire, H5 évalue les places vides. Grâce à un mécanisme de "Look-Ahead", l'algorithme lit les cibles futures de tous les autres bacs. S'il s'apprête à déposer le bac bloquant sur une case vide qui s'avère être la \textbf{cible finale} d'un autre bac, il applique une forte pénalité mathématique (ajoute 100 à la distance fictive). Cela évite de "polluer" inutilement une case précieuse.

\subsection{H6 : Macro-Structurelle / RTLM (Linear Merge)}
S'éloignant du rangement ciblé, H6 modifie fondamentalement l'ordre de traitement. Elle trie les bacs à traiter en se basant sur la \textbf{colonne de leur destination finale}. Le solveur tente ainsi de construire et de ranger l'entrepôt "colonne par colonne", comme on construirait un mur brique par brique.

\subsection{H7 : Topologie Respirante}
Une stratégie novatrice axée sur le désengorgement de la zone de travail. Lors de la phase de dégagement (Double Handling), au lieu de chercher la case vide la plus proche (ce qui maintient le bac bloquant dans les pattes de la navette), H7 envoie délibérément le bac pénalisant \textbf{au fond absolu de l'entrepôt} (la colonne ayant l'index le plus élevé possible). Cela permet à la zone dense de "respirer" et offre un champ d'action dégagé pour ranger correctement le reste des composants.

\section{L'Outil de Benchmarking (Évaluation)}
Le script embarque la fonction \texttt{evaluer\_heuristiques}. Elle agit comme un laboratoire automatisé. Fournie avec un dossier d'instances générées (par \texttt{instance\_generator.py}), elle instancie un \texttt{Solveur} pour chaque fichier, fait tourner les 7 heuristiques séquentiellement, et centralise les métriques analytiques (Mouvements totaux, Taux de complétion, Temps CPU, Quantité de Double Handling) dans un rapport CSV exhaustif, prêt à être exploité statistiquement.

\end{document}
